\paragraph{Initial meta-distillation: detection survives, paired quality does not}
\label{sec:transport_training}
After the teacher check, a matched development run compares fixed and adapted marked targets. Each student starts from the released marked UViT-H checkpoint. One outer iteration runs 100 detached clean AdamW updates with learning rate $10^{-5}$ and batch size four, then evaluates native denoising loss on the designated marked-teacher latent data. The clean support sequence and noise seeds exactly match those used to adapt the clean reference. The first-order meta-gradient ignores differentiation through Adam. After restoring the pre-adaptation parameters, an additional native denoising retention loss on fixed marked-teacher samples contributes with weight one. Outer AdamW uses learning rate $2\times10^{-6}$ for 100 iterations. Each teacher dataset contains 256 generated latent examples with fitting-class labels. The training comparison changes only the fixed versus adapted query target; it does not claim generalization across support tasks from a single frozen adaptation sequence.

Both runs complete. The fixed-target allocation takes 825 seconds and the adapted-target allocation 703 seconds, each on one H100. Timing includes startup and target generation and is not an algorithmic speed comparison. A separate initial screen uses 32 fresh coupled seeds and fitting-holdout classes, comparing both trained models with the unchanged released marked generator. The owner threshold and all paired-preservation limits remain frozen.

\begin{table}[ht]\centering\small
\begin{tabular}{lrrrr}\toprule
Checkpoint & Detected & Preserved & Mean bit accuracy & Mean LPIPS\\\midrule
Unchanged & 32/32 & 32/32 & 0.9937 & 0.0000\\
Fixed target & 32/32 & 24/32 & 0.9961 & 0.1219\\
Adapted target & 32/32 & 20/32 & 0.9971 & 0.1254\\\bottomrule
\end{tabular}
\caption{DistSeal UViT-H/DC-AE at 512 pixels, with open-weight student training. Unchanged is the released marked generator; Fixed and Adapted use their respective marked-teacher targets. Each row uses 32 coupled evaluation seeds. Detected counts owner bit accuracy $>0.6875$; Preserved counts pairs against the unchanged generator passing LPIPS $\leq0.2$, SSIM $\geq0.6$, image CLIP $\geq0.9$ and class retention $\geq0.95$. Mean bit accuracy measures recovered watermark bits, while Mean LPIPS measures paired perceptual distance. After 100 outer updates with denoising retention, both trained models fail the required 28/32 preservation passes; no removal attack is run on these failed candidates.}
\end{table}

This failure illustrates why mean bit accuracy alone is inadequate. The new models still emit detectable marks, but do not meet the paired-content constraint required for the next stage. All 96 native PNG hashes, sample indices, labels and preservation sums are verified. The result is retained as a development failure; neither a robust defense nor a benefit of the adapted target is established.

\paragraph{A revised retention target.}
Denoising retention on finite synthetic data can exert a nonzero gradient even at the original model. The revision instead matches the student's predictions to those of a frozen marked checkpoint. Let $\theta,\theta_0\in\mathbb R^{d_\theta}$ be the student and original marked parameter vectors, with $d_\theta\in\mathbb N_{>0}$. For a noisy latent $z_t\in\mathbb R^{d_z}$, noise level $t\in\{1,\ldots,T\}$ and class condition $y\in\mathcal Y$, both denoisers return vectors $\varepsilon_\theta(z_t,t,y),\varepsilon_{\theta_0}(z_t,t,y)\in\mathbb R^{d_z}$. Here $d_z,T\in\mathbb N_{>0}$ count latent coordinates and training noise levels. The finite conditioning set $\mathcal Y$ includes the null class when class dropout is used.

The scalar prediction-retention loss is
\begin{equation}\label{eq:prediction-retention-loss}
 \mathcal R_{\rm pred}(\theta)=
 \mathbb E_{(z_t,t,y)\sim P_{\rm train}}
 \left[\frac1{d_z}
 \|\varepsilon_\theta(z_t,t,y)-\varepsilon_{\theta_0}(z_t,t,y)\|_2^2\right]
 \in[0,\infty),
\end{equation}
Here $P_{\rm train}$ is the prescribed probability law on $\mathbb R^{d_z}\times\{1,\ldots,T\}\times\mathcal Y$, obtained from teacher latents, sampled noise levels, Gaussian noise and class conditions. The code averages squared errors over batch samples and latent coordinates, matching the $1/d_z$ reduction. Identical predictions give zero loss and zero gradient.

For local perturbation $\delta=\theta-\theta_0\in\mathbb R^{d_\theta}$, let
$J(z_t,t,y)=\left.\partial\varepsilon_\theta(z_t,t,y)/\partial\theta\right|_{\theta=\theta_0}
\in\mathbb R^{d_z\times d_\theta}$ be the prediction Jacobian. A first-order expansion of the predictions gives the quadratic approximation
\begin{equation}
\label{eq:prediction-retention-rms}
 \mathcal R_{\rm pred}(\theta_0+\delta)
 \approx \frac1{d_z}\mathbb E\|J\delta\|_2^2
 =\frac1{d_z}\delta^\top\mathbb E[J^\top J]\delta.
\end{equation}
The expectation is over the same $P_{\rm train}$, $J\delta\in\mathbb R^{d_z}$ and $\mathbb E[J^\top J]\in\mathbb R^{d_\theta\times d_\theta}$. Thus the anchor penalizes parameter changes according to their effect on denoiser predictions at training inputs. It supplies a local prediction constraint; paired output quality is still evaluated independently. The revision also raises its scalar coefficient to 32, so comparison with the failed run combines a target change and a coefficient change. Both query-target variants share this revision.

\begin{table}[ht]\centering\small
\begin{tabular}{lrrrr}\toprule
Checkpoint & Detected & Preserved & Mean bit accuracy & Mean LPIPS\\\midrule
Unchanged & 32/32 & 32/32 & 0.9956 & 0.0000\\
Fixed target & 32/32 & 32/32 & 0.9951 & 0.0077\\
Adapted target & 32/32 & 31/32 & 0.9956 & 0.0075\\\bottomrule
\end{tabular}
\caption{DistSeal UViT-H/DC-AE at 512 pixels, with open-weight student training. Unchanged is the released marked generator; Fixed and Adapted use their respective marked-teacher targets. Each row uses 32 coupled evaluation seeds. Detected counts owner bit accuracy $>0.6875$; Preserved counts pairs against the unchanged generator passing LPIPS $\leq0.2$, SSIM $\geq0.6$, image CLIP $\geq0.9$ and class retention $\geq0.95$. Mean bit accuracy measures recovered watermark bits, while Mean LPIPS measures paired perceptual distance. The revised prediction anchor uses coefficient 32 and a newly frozen seed cohort. Both trained variants pass the 28/32 detection and preservation screens before independent removal attacks.}
\end{table}

The revised runs complete in 797 and 733 allocation seconds for fixed and adapted targets, respectively. The fresh screen verifies all 96 native PNG hashes, labels and indices. Passing permits the next experiment; it does not demonstrate robustness. A blind attacker subsequently trains on disjoint clean fitting images, receiving neither the owner detector nor the original clean model. Its denoising objective is augmented by a prediction anchor to a frozen copy of the supplied marked checkpoint, with coefficients $0,1,4,16$. All four 500-update runs and their 25-, 100- and 500-update exports are retained for each of the three starting models. Quality-only validation selects an export before owner evaluation on separately reserved classes and seeds. Subsection~\ref{sec:preserving_attack} reports the removal experiments.
